\documentclass[aspectratio=169, xcolor=dvipsnames]{beamer}

% --- Theme Setup ---
\usetheme{Madrid}
\usecolortheme{default}

% --- Packages ---
\usepackage{graphicx}
\usepackage{xcolor}
\usepackage{booktabs}
\usepackage{hyperref}
\usepackage{adjustbox}
\usepackage{amssymb}
\usepackage{amsmath}
\usepackage{listings}
\usepackage{caption}

\lstset{
  basicstyle=\ttfamily\small,
  keywordstyle=\color{blue}\bfseries,
  stringstyle=\color{red},
  showstringspaces=false,
  breaklines=true
}

% --- Title Information ---
\title{Module 14}
\subtitle{Intermediate SQL/3}
\author{Milav Dabgar}
\institute{www.milav.in}
\date{\today}

\begin{document}

% Title Slide
\begin{frame}
    \titlepage
\end{frame}

\begin{frame}{Module Recap}
    \begin{itemize}
        \item SQL expressions for Join and Views
    \end{itemize}
\end{frame}

\begin{frame}{Module Objectives}
    \begin{itemize}
        \item To understand Transactions
        \item To learn SQL expressions for Integrity Constraints
        \item To understand more Data Types in SQL
        \item To understand Authorization in SQL
    \end{itemize}
\end{frame}

\begin{frame}{Module Outline}
    \begin{itemize}
        \item Transactions
        \item Integrity Constraints
        \item SQL Data Types and Schemas
        \item Authorization
    \end{itemize}
\end{frame}

\section{Transactions}

\begin{frame}
  \vfill
  \centering
  \begin{beamercolorbox}[sep=8pt,center,shadow=true,rounded=true]{title}
    \usebeamerfont{title}Transactions\par%
  \end{beamercolorbox}
  \vfill
\end{frame}

\begin{frame}{Transactions}
    \begin{itemize}
        \item Unit of work
        \item \textbf{Atomic transaction}
        \begin{itemize}
            \item either fully executed or rolled back as if it never occurred
        \end{itemize}
        \item Isolation from concurrent transactions
        \item Transactions begin implicitly
        \item Ended by \texttt{commit work} or \texttt{rollback work}
        \item But default on most databases: each SQL statement commits automatically
        \item Can turn off auto commit for a session (for example, using API)
        \item In SQL:1999, can use: \texttt{begin atomic ... end}
        \begin{itemize}
            \item Not supported on most databases
        \end{itemize}
    \end{itemize}
\end{frame}

\section{Integrity Constraints}

\begin{frame}
  \vfill
  \centering
  \begin{beamercolorbox}[sep=8pt,center,shadow=true,rounded=true]{title}
    \usebeamerfont{title}Integrity Constraints\par%
  \end{beamercolorbox}
  \vfill
\end{frame}

\begin{frame}{Integrity Constraints}
    \begin{itemize}
        \item Integrity constraints guard against accidental damage to the database, by ensuring that authorized changes to the database do not result in a loss of data consistency
        \item A checking account must have a balance greater than Rs. 10,000.00
        \item A salary of a bank employee must be at least Rs. 250.00 an hour
        \item A customer must have a (non-null) phone number
    \end{itemize}
\end{frame}

\begin{frame}{Integrity Constraints on a Single Relation}
    \begin{itemize}
        \item \texttt{not null}
        \item \texttt{primary key}
        \item \texttt{unique}
        \item \texttt{check (P)}, where \texttt{P} is a predicate
    \end{itemize}
\end{frame}

\begin{frame}[fragile]{Not Null and Unique Constraints}
    \begin{itemize}
        \item \textbf{\texttt{not null}}
        \begin{itemize}
            \item Declare name and budget to be \texttt{not null}
            \begin{lstlisting}[language=SQL]
name varchar(20) not null 
budget numeric(12,2) not null
            \end{lstlisting}
        \end{itemize}
        \item \textbf{\texttt{unique ($A_1, A_2, ..., A_m$)}}
        \begin{itemize}
            \item The \texttt{unique} specification states that the attributes $A_1, A_2, ..., A_m$ form a candidate key
            \item Candidate keys are permitted to be null (in contrast to primary keys).
        \end{itemize}
    \end{itemize}
\end{frame}

\begin{frame}[fragile]{The check clause}
    \begin{itemize}
        \item \texttt{check (P)}, where \texttt{P} is a predicate
        \item Ensure that semester is one of fall, winter, spring or summer:
        \begin{lstlisting}[language=SQL, basicstyle=\ttfamily\scriptsize]
create table section ( 
    course_id varchar(8), 
    sec_id varchar(8), 
    semester varchar(6), 
    year numeric(4,0), 
    building varchar(15), 
    room_number varchar(7), 
    time_slot_id varchar(4), 
    primary key (course_id, sec_id, semester, year), 
    check (semester in ('Fall', 'Winter', 'Spring', 'Summer')) 
);
        \end{lstlisting}
    \end{itemize}
\end{frame}

\section{Referential Integrity}

\begin{frame}{Referential Integrity}
    \begin{itemize}
        \item Ensures that a value that appears in one relation for a given set of attributes also appears for a certain set of attributes in another relation
        \item Example: If 'Biology' is a department name appearing in one of the tuples in the \texttt{instructor} relation, then there exists a tuple in the \texttt{department} relation for 'Biology'
        \item Let $A$ be a set of attributes. Let $R$ and $S$ be two relations that contain attributes $A$ and where $A$ is the primary key of $S$. $A$ is said to be a \textbf{foreign key} of $R$ if for any values of $A$ appearing in $R$ these values also appear in $S$.
    \end{itemize}
\end{frame}

\begin{frame}[fragile]{Cascading Actions in Referential Integrity}
    \begin{itemize}
        \item With cascading, you can define the actions that the Database Engine takes when a user tries to delete or update a key to which existing foreign keys point
        \begin{lstlisting}[language=SQL, basicstyle=\ttfamily\scriptsize]
create table course (
    course_id char(5) primary key, 
    title varchar(20), 
    dept_name varchar(20) references department
);
        \end{lstlisting}
        \vfill
        \begin{lstlisting}[language=SQL, basicstyle=\ttfamily\scriptsize]
create table course (
    ...
    dept_name varchar(20), 
    foreign key (dept_name) references department 
        on delete cascade 
        on update cascade,
    ...
);
        \end{lstlisting}
        \item Alternative actions to cascade: \texttt{no action}, \texttt{set null}, \texttt{set default}
    \end{itemize}
\end{frame}

\begin{frame}[fragile]{Integrity Constraint Violation During Transactions}
    \begin{itemize}
        \item \begin{lstlisting}[language=SQL, basicstyle=\ttfamily\scriptsize]
create table person ( 
    ID char(10), 
    name char(40), 
    mother char(10), 
    father char(10), 
    primary key ID, 
    foreign key father references person, 
    foreign key mother references person
);
        \end{lstlisting}
        \item How to insert a tuple without causing constraint violation?
        \begin{itemize}
            \item Insert father and mother of a person before inserting person
            \item OR, Set father and mother to null initially, update after inserting all persons (not possible if father and mother attributes declared to be \texttt{not null})
            \item OR Defer constraint checking (will discuss later)
        \end{itemize}
    \end{itemize}
\end{frame}

\section{SQL Data Types and Schemas}

\begin{frame}
  \vfill
  \centering
  \begin{beamercolorbox}[sep=8pt,center,shadow=true,rounded=true]{title}
    \usebeamerfont{title}SQL Data Types and Schemas\par%
  \end{beamercolorbox}
  \vfill
\end{frame}

\begin{frame}{Built-in Data Types in SQL}
    \begin{itemize}
        \item \texttt{date}: Dates, containing a (4 digit) year, month and date
        \begin{itemize}
            \item Example: \texttt{date '2005-7-27'}
        \end{itemize}
        \item \texttt{time}: Time of day, in hours, minutes and seconds.
        \begin{itemize}
            \item Example: \texttt{time '09:00:30'}, \texttt{time '09:00:30.75'}
        \end{itemize}
        \item \texttt{timestamp}: date plus time of day
        \begin{itemize}
            \item Example: \texttt{timestamp '2005-7-27 09:00:30.75'}
        \end{itemize}
        \item \texttt{interval}: period of time
        \begin{itemize}
            \item Example: \texttt{interval '1' day}
        \end{itemize}
        \item Subtracting a date/time/timestamp value from another gives an interval value
        \item Interval values can be added to date/time/timestamp values
    \end{itemize}
\end{frame}

\begin{frame}[fragile]{Index Creation}
    \begin{itemize}
        \item \begin{lstlisting}[language=SQL, basicstyle=\ttfamily\scriptsize]
create table student ( 
    ID varchar(5), 
    name varchar(20) not null, 
    dept_name varchar(20), 
    tot_cred numeric(3,0) default 0, 
    primary key (ID)
);
        \end{lstlisting}
        \item \begin{lstlisting}[language=SQL, basicstyle=\ttfamily\scriptsize]
create index studentID_index on student(ID);
        \end{lstlisting}
        \item Indices are data structures used to speed up access to records with specified values for index attributes
        \item \begin{lstlisting}[language=SQL, basicstyle=\ttfamily\scriptsize]
select * 
from student 
where ID = '12345';
        \end{lstlisting}
        \item Can be executed by using the index to find the required record, without looking at all records of \texttt{student}
        \item More on indices in Chapter 9
    \end{itemize}
\end{frame}

\begin{frame}[fragile]{User-Defined Types}
    \begin{itemize}
        \item \texttt{create type} construct in SQL creates user-defined type (alias, like \texttt{typedef} in C)
        \item \begin{lstlisting}[language=SQL]
create type Dollars as numeric(12,2) final;
        \end{lstlisting}
        \item \begin{lstlisting}[language=SQL]
create table department (
    dept_name varchar(20), 
    building varchar(15), 
    budget Dollars 
);
        \end{lstlisting}
    \end{itemize}
\end{frame}

\begin{frame}[fragile]{Domains}
    \begin{itemize}
        \item \texttt{create domain} construct in SQL-92 creates user-defined domain types 
        \begin{lstlisting}[language=SQL]
create domain person_name char(20) not null;
        \end{lstlisting}
        \item Types and domains are similar
        \item Domains can have constraints, such as \texttt{not null}, specified on them
        \begin{lstlisting}[language=SQL, basicstyle=\ttfamily\scriptsize]
create domain degree_level varchar(10) 
    constraint degree_level_test 
    check (value in ('Bachelors', 'Masters', 'Doctorate'));
        \end{lstlisting}
    \end{itemize}
\end{frame}

\begin{frame}{Large-Object Types}
    \begin{itemize}
        \item Large objects (photos, videos, CAD files, etc.) are stored as a \textbf{large object}:
        \begin{itemize}
            \item \textbf{\texttt{blob}}: binary large object - object is a large collection of uninterpreted binary data (whose interpretation is left to an application outside of the database system)
            \item \textbf{\texttt{clob}}: character large object - object is a large collection of character data
        \end{itemize}
        \item When a query returns a large object, a pointer is returned rather than the large object itself
    \end{itemize}
\end{frame}

\section{Authorization}
\begin{frame}
  \vfill
  \centering
  \begin{beamercolorbox}[sep=8pt,center,shadow=true,rounded=true]{title}
    \usebeamerfont{title}Authorization\par%
  \end{beamercolorbox}
  \vfill
\end{frame}

\begin{frame}{Authorization}
    \begin{itemize}
        \item Forms of authorization on parts of the database:
        \begin{itemize}
            \item \textbf{Read} - allows reading, but not modification of data
            \item \textbf{Insert} - allows insertion of new data, but not modification of existing data
            \item \textbf{Update} - allows modification, but not deletion of data
            \item \textbf{Delete} - allows deletion of data
        \end{itemize}
        \item Forms of authorization to modify the database schema
        \begin{itemize}
            \item \textbf{Index} - allows creation and deletion of indices
            \item \textbf{Resources} - allows creation of new relations
            \item \textbf{Alteration} - allows addition or deletion of attributes in a relation
            \item \textbf{Drop} - allows deletion of relations
        \end{itemize}
    \end{itemize}
\end{frame}

\begin{frame}{Authorization Specification in SQL}
    \begin{itemize}
        \item The \texttt{grant} statement is used to confer authorization
        \begin{center}
        \texttt{grant <privilege list> on <relation name or view name> to <user list>}
        \end{center}
        \item \texttt{<user list>} is:
        \begin{itemize}
            \item a user-id
            \item \texttt{public}, which allows all valid users the privilege granted
            \item A \textbf{role} (more on this later)
        \end{itemize}
        \item Granting a privilege on a view does not imply granting any privileges on the underlying relations
        \item The grantor of the privilege must already hold the privilege on the specified item (or be the database administrator)
    \end{itemize}
\end{frame}

\begin{frame}[fragile]{Privileges in SQL}
    \begin{itemize}
        \item \textbf{\texttt{select}}: allows read access to relation, or the ability to query using the view
        \begin{itemize}
            \item Example: grant users $U_1, U_2,$ and $U_3$ select authorization on the \texttt{instructor} relation:
            \begin{lstlisting}[language=SQL]
grant select on instructor to U1, U2, U3;
            \end{lstlisting}
        \end{itemize}
        \item \textbf{\texttt{insert}}: the ability to insert tuples
        \item \textbf{\texttt{update}}: the ability to update using the SQL update statement
        \item \textbf{\texttt{delete}}: the ability to delete tuples
        \item \textbf{\texttt{all privileges}}: used as a short form for all the allowable privileges
    \end{itemize}
\end{frame}

\begin{frame}[fragile]{Revoking Authorization in SQL}
    \begin{itemize}
        \item The \texttt{revoke} statement is used to revoke authorization
        \begin{center}
        \texttt{revoke <privilege list> on <relation name or view name> from <user list>}
        \end{center}
        \item Example:
        \begin{lstlisting}[language=SQL]
revoke select on branch from U1, U2, U3;
        \end{lstlisting}
        \item \texttt{<privilege-list>} may be \texttt{all} to revoke all privileges the revokee may hold
        \item If \texttt{<revokee-list>} includes \texttt{public}, all users lose the privilege except those granted it explicitly
        \item If the same privilege was granted twice to the same user by different grantees, the user may retain the privilege after the revocation
        \item All privileges that depend on the privilege being revoked are also revoked
    \end{itemize}
\end{frame}

\begin{frame}[fragile]{Roles}
    \begin{itemize}
        \item \begin{lstlisting}[language=SQL]
create role instructor;
grant instructor to Amit;
        \end{lstlisting}
        \item Privileges can be granted to roles:
        \begin{lstlisting}[language=SQL]
grant select on takes to instructor;
        \end{lstlisting}
        \item Roles can be granted to users, as well as to other roles
        \begin{lstlisting}[language=SQL]
create role teaching_assistant;
grant teaching_assistant to instructor;
        \end{lstlisting}
        \item Instructor inherits all privileges of teaching assistant
        \item \textbf{Chain of roles}
        \begin{lstlisting}[language=SQL]
create role dean;
grant instructor to dean;
grant dean to Satoshi;
        \end{lstlisting}
    \end{itemize}
\end{frame}

\begin{frame}[fragile]{Authorization on Views}
    \begin{itemize}
        \item \begin{lstlisting}[language=SQL]
create view geo_instructor as (
    select * 
    from instructor 
    where dept_name = 'Geology'
);
grant select on geo_instructor to geo_staff;
        \end{lstlisting}
        \item Suppose that a \texttt{geo\_staff} member issues:
        \begin{lstlisting}[language=SQL]
select * from geo_instructor;
        \end{lstlisting}
        \item What if
        \begin{itemize}
            \item \texttt{geo\_staff} does not have permissions on \texttt{instructor}?
            \item creator of view did not have some permissions on \texttt{instructor}?
        \end{itemize}
    \end{itemize}
\end{frame}

\begin{frame}[fragile]{Other Authorization Features}
    \begin{itemize}
        \item \textbf{\texttt{references}} privilege to create foreign key
        \begin{lstlisting}[language=SQL]
grant reference (dept_name) on department to Mariano;
        \end{lstlisting}
        \item why is this required?
        \item \textbf{Transfer of privileges}
        \begin{lstlisting}[language=SQL]
grant select on department to Amit with grant option;
revoke select on department from Amit, Satoshi cascade;
revoke select on department from Amit, Satoshi restrict;
        \end{lstlisting}
    \end{itemize}
\end{frame}

\begin{frame}{Module Summary}
    \begin{itemize}
        \item Introduced transactions
        \item Learnt SQL expressions for integrity constraints
        \item Familiarized with more data types in SQL
        \item Discussed authorization in SQL
    \end{itemize}
    \vspace{2em}
    \begin{center}
    \small \textit{Slides used in this presentation are borrowed from \url{http://db-book.com/} with kind permission of the authors.} \\
    \small \textit{Edited and new slides are marked with 'PPD'.}
    \end{center}
\end{frame}

\end{document}
